% =============================================================================
%  Foundations of Computational Theology:
%  Set-Theoretic Meta-Ethics and the Gödelian Completion
%
%  Author: Rafael D. De Paz
%  Date:   March 2026
%  Target: Theology and Science / Zygon
% =============================================================================
\documentclass[12pt, a4paper]{article}
\usepackage[utf8]{inputenc}
\usepackage[T1]{fontenc}
\usepackage{amsmath, amssymb, amsthm, mathtools}
\usepackage{geometry}
\usepackage{hyperref}
\usepackage{cleveref}
\usepackage{enumitem}
\usepackage{booktabs}
\usepackage{microtype}
\usepackage{natbib}
\geometry{margin=1in}

\theoremstyle{plain}
\newtheorem{theorem}{Theorem}[section]
\newtheorem{lemma}[theorem]{Lemma}
\newtheorem{proposition}[theorem]{Proposition}
\newtheorem{corollary}[theorem]{Corollary}
\theoremstyle{definition}
\newtheorem{definition}[theorem]{Definition}
\newtheorem{axiom}[theorem]{Axiom}
\theoremstyle{remark}
\newtheorem{remark}[theorem]{Remark}

\hypersetup{
  colorlinks=true, linkcolor=blue, citecolor=blue, urlcolor=blue,
  pdftitle={Foundations of Computational Theology},
  pdfauthor={Rafael D. De Paz}
}

\begin{document}

\title{\textbf{Foundations of Computational Theology: \\
Set-Theoretic Meta-Ethics and the \\
G\"odelian Completion}}
\author{Rafael D. De Paz \\ \textit{Independent Researcher} \\ \texttt{hi@logos.pub}}
\date{March 2026}
\maketitle

\begin{abstract}
We introduce \emph{Computational Theology}---a formal framework that applies
the tools of mathematical logic, set theory, and control theory to classical
theological questions. We present three foundational results. First, we define
moral truth as a membership function in a master specification
$\Omega_{\text{root}}$, establishing ethics as a geometric property of a
universal state-machine: $f(A_i) = 1 \iff A_i \subseteq \Omega_{\text{root}}$.
Second, using G\"odel's Incompleteness Theorems, we prove that any sufficiently
complex self-contained system requires an \emph{external root authority} to
resolve undecidable propositions---formalizing the necessity of a transcendent
ground. Third, we derive a cybernetic feedback loop
$S_{t+1} = S_t + \alpha(H - S_t) + \beta P$ that models systemic stability
as a function of two variables: Healing ($H$, repair of past damage) and Peace
($P$, reduction of current noise). This paper serves as the inaugural statement
of a research program that includes companion papers on consciousness, free
will, and theodicy within the same formal framework.
\end{abstract}

\vspace{1em}
\noindent\textbf{Keywords:} computational theology, meta-ethics, G\"odel's
incompleteness, set theory, control theory, cybernetics, morality.

\section{Introduction}\label{sec:intro}

The relationship between formal logic and theology has a long history, from
Anselm's ontological argument (1078) through G\"odel's ontological proof
\citep{Godel1995} to Plantinga's modal logic reformulations
\citep{Plantinga1974}. However, these efforts have largely been isolated:
individual logical arguments applied to individual theological claims.

We propose a unified framework---\emph{Computational Theology}---that
treats theology as a \emph{formal system} subject to the same constraints
as any other: consistency, completeness, decidability, and entropy. This
is not a metaphor. We claim that theological propositions can be formalized,
that ethical truths have geometric structure, and that the necessity of a
transcendent ground can be \emph{derived} from G\"odel's theorems.

\subsection{Main Results}
\begin{enumerate}
  \item The \textbf{Alignment Function}: Moral truth as set membership
    in $\Omega_{\text{root}}$ (\S\ref{sec:alignment}).
  \item The \textbf{G\"odelian Completion}: Necessity of an external root
    authority (\S\ref{sec:godelian}).
  \item The \textbf{Cybernetic Stability Loop}: Systemic convergence under
    $S = H + P$ (\S\ref{sec:cybernetic}).
  \item The \textbf{Error Taxonomy}: Classification of ethical failure modes
    as computational errors (\S\ref{sec:errors}).
\end{enumerate}

\section{The Alignment Function}\label{sec:alignment}

\begin{definition}[Master Specification]\label{def:omega}
  Let $\Omega_{\text{root}}$ be a consistent, complete specification of the
  optimal state of a universal system. An action $A_i$ is \emph{aligned}
  (morally correct) if and only if:
  \begin{equation}\label{eq:alignment}
    f(A_i) = 1 \iff A_i \subseteq \Omega_{\text{root}}.
  \end{equation}
\end{definition}

\begin{axiom}[Geometric Morality]\label{ax:geometric}
  Moral truth is a rigid geometric property of the universal state-machine,
  not a subjective preference. The alignment function $f$ is deterministic
  and binary: an action either conforms to $\Omega_{\text{root}}$ or it
  does not.
\end{axiom}

\begin{definition}[Computational Sin]\label{def:sin}
  An action $A_i$ with $f(A_i) = 0$ is a \emph{computational error}---it
  attempts to execute outside the allocated address space of
  $\Omega_{\text{root}}$. We classify errors as:
  \begin{enumerate}[label=(\roman*)]
    \item \textbf{Memory Leak:} Resources consumed without productive output.
    \item \textbf{Segmentation Fault:} Access to forbidden address space.
    \item \textbf{Stack Overflow:} Recursive self-reference without termination.
  \end{enumerate}
\end{definition}

\section{The G\"odelian Completion}\label{sec:godelian}

\begin{theorem}[Necessity of External Root]\label{thm:godelian}
  Let $\mathcal{S}$ be any formal system sufficiently powerful to encode
  arithmetic (and therefore physics, ethics, or any domain with comparable
  expressive power). By G\"odel's First Incompleteness Theorem
  \citep{Godel1931}:
  \begin{enumerate}[label=(\alph*)]
    \item $\mathcal{S}$ contains true statements that cannot be proved
      within $\mathcal{S}$.
    \item The consistency of $\mathcal{S}$ cannot be established from
      within $\mathcal{S}$ (Second Incompleteness Theorem).
  \end{enumerate}
  Therefore, $\mathcal{S}$ requires an \emph{external oracle}---a root
  authority not contained within $\mathcal{S}$---to resolve undecidable
  propositions and verify its own consistency.
\end{theorem}

\begin{remark}[The Theological Mapping]
  This is a formal restatement of the classical cosmological argument:
  a self-contained system cannot ground its own existence. The external
  oracle in G\"odel's framework corresponds structurally to the concept
  of a transcendent ground in classical theology.
\end{remark}

\section{The Cybernetic Stability Loop}\label{sec:cybernetic}

\begin{definition}[System Stability]\label{def:stability}
  The stability $S$ of a system at time $t$ evolves according to:
  \begin{equation}\label{eq:stability}
    S_{t+1} = S_t + \alpha(H - S_t) + \beta P,
  \end{equation}
  where $H$ is the Healing function (repair of accumulated damage),
  $P$ is the Peace function (reduction of current environmental noise),
  and $\alpha, \beta > 0$ are gain coefficients representing the agent's
  ``will'' to execute repair.
\end{definition}

\begin{theorem}[Convergence to Stability]\label{thm:convergence}
  If $H > 0$ and $P > 0$ remain non-zero, the system asymptotically
  converges: $\lim_{t \to \infty} S_t = S^*$, where $S^*$ is the fixed
  point satisfying $S^* = \alpha H / (1 + \alpha - 1) + \beta P$.
  If $H = 0$ (no healing), the system drifts into oscillatory chaos.
\end{theorem}

\section{Error Taxonomy}\label{sec:errors}

\begin{center}
\begin{tabular}{lll}
  \toprule
  \textbf{Computational Error} & \textbf{Ethical Mapping} & \textbf{Correction} \\
  \midrule
  Memory Leak & Acedia / Sloth & Garbage collection (repentance) \\
  Segmentation Fault & Transgression & Access control (conscience) \\
  Stack Overflow & Obsession / Addiction & Recursion limit (humility) \\
  Race Condition & Injustice & Mutex lock (justice protocols) \\
  \bottomrule
\end{tabular}
\end{center}

\section{Discussion}\label{sec:discussion}

This paper inaugurates a research program in Computational Theology. Companion
papers develop specific applications:
\begin{enumerate}
  \item \textbf{Consciousness as Recursive Self-Reference} --- Formalizes
    the observer as a G\"odelian oracle.
  \item \textbf{Free Will as Computational Necessity} --- Proves agency
    is required for system novelty.
  \item \textbf{Theodicy as Error Handling} --- Reframes the Problem of
    Evil as fault tolerance in distributed systems.
\end{enumerate}

\subsection{Limitations}
\begin{enumerate}
  \item The alignment function $f$ assumes the existence and consistency
    of $\Omega_{\text{root}}$---this is an axiom, not a derivation.
  \item The G\"odelian argument establishes \emph{structural necessity}
    of an external ground, not its specific nature or attributes.
  \item The cybernetic model is a first approximation; real ethical
    systems involve higher-order feedback loops.
\end{enumerate}

\section{Conclusion}\label{sec:conclusion}

Computational Theology provides a formal language for theological reasoning
that is consistent, falsifiable within its axioms, and interoperable with
the existing tools of mathematics and computer science. The framework
demonstrates that ethical truths have geometric structure, that self-contained
systems require external grounding, and that systemic stability is achievable
through the sustained application of healing and peace.

\bibliographystyle{plainnat}
\bibliography{references}

\end{document}
